\section{Gaussian geometry before the first redundancy}\label{sec:geometry}

\begin{lemma}[Transverse Gaussian bound]\label{lem:smallball}
Let $(D,U)$ be jointly Gaussian in $\R^3\times\R^3$, with $U$ centered
and isotropic. If $\lambda=\lambda_{\max}\Cov(D\mid U)>0$, then
\[
 \Pp\{\operatorname{dist}(D,\R U)\leq\eps\}
       \leq C\min(1,\eps/\sqrt\lambda).
\]
\end{lemma}
\begin{proof}
Choose a top eigenvector $e$. Conditional on $U=u$, project onto
$w=P_{u^\perp}e/\sin\theta$, where $\theta$ is the angle of $u,e$.
The variance of $w\cdot D$ is at least $\lambda\sin^2\theta$.
A one-dimensional Gaussian density gives
$C\min(1,\eps/(\sqrt\lambda\sin\theta))$, independently of its mean.
The direction of $U$ is uniform and $\E(\sin\theta)^{-1}=\pi/2$.
The zero and parallel cases are null.
\end{proof}

\subsection{Repeated earlier pairs}
\begin{lemma}[An erased normal coordinate]\label{lem:repeated}
Fix a chronological merger forest, its normals and C/O flags.
Suppose A and B have a prescribed contact at $t_0$, and the prospective
first redundant contact repeats this pair. Let $t_1>t_0$ be the first
later C contact touching either endpoint, and $\ell=t_1-t_0$.
At the beginning of the tested subsequent free segment, write $D$ for
their relative displacement and $Y$ for both endpoint velocities. Then
\[
 \Cov(D\mid Y)\succeq\sigma^2\ell^2 n_1n_1^T,
 \qquad
 \Pp_\gamma\{\text{return}\}\leq C_\sigma\min(1,\eps/\ell).
\]
If no such C contact exists, a second incoming crossing is impossible.
\end{lemma}
\begin{proof}
A C collision swaps the two normal components and retains the transverse
ones. It is orthogonal in all birth Gaussian coordinates; the current
velocity array is therefore a product Gaussian at fixed free normals.
After A and B first belong to one component, every merger touching an
endpoint imports a disjoint component. Its incident normal scalar is
Gaussian and independent of the entire previous tagged history.
That imported component cannot feed a scalar back through another merger:
such a contact would already be redundant.

Suppose A has the first C update. With $A_0,B_0$ the velocities just after
the earlier pair contact, set $\xi=A_0\cdot n_1$. It is independent of
$P_{n_1^\perp}A_0,B_0$ and all later imported scalars. At $t_1$ it is
erased from A and never re-enters a tagged endpoint before the first
redundancy. Consequently
\[
 D=\ell n_1\xi+R,\qquad
 \xi\ \text{is independent of }(R,Y).
\]
The initial contact displacement is a deterministic $\eps$-normal
offset. This proves the covariance inequality. Updating B changes only
the sign. The final relative velocity is isotropic, and the proof of
\Cref{lem:smallball} gives the stated probability bound.
Without a C update the relative line remains straight: a C contact
separates, while an O contact has only the outgoing exit. Neither permits
a second incoming crossing.
\end{proof}

The conditioning is only on the two endpoint velocities, not on all
partner velocities, realizability, or collision fluxes. The erased scalar
was exported to a partner and can be revealed by those stronger
conditionings.

\subsection{Nonneighbor pairs}
\begin{lemma}[Rank one across a forest cut]\label{lem:rankone}
The $3$ by $3$ velocity-transfer block between distinct labels in a
prescribed forest of C/O contacts has rank at most one.
\end{lemma}
\begin{proof}
Cut an edge on their unique path. It is the only contact crossing that
cut. A C contact transmits one normal scalar across it; O transmits none.
Every later update is internal to one side, so dependence on the original
velocity continues to factor through that scalar. This does not require
the chronological order to be monotone along the spatial path.
\end{proof}

Fix nonneighbors A,B. Orient tree edges $e=(p_e,q_e)$, put
$b_e=e_{p_e}-e_{q_e}$, and write
$b=e_A-e_B=\sum_es_eb_e$ for their signed spatial path.
The actual contact equations give, at a time $t_*$ after the prefix,
\begin{equation}\label{eq:position}
 D_\eps=\eps\sum_es_e\omega_e+
   \sum_es_e\int_{t_e}^{t_*}(v_{p_e}(r)-v_{q_e}(r))\dd r .
\end{equation}
This telescopes positions; it makes no constrained initial positions
independent. At fixed normals all velocities are linear in the birth array
$V$. Write $D_\eps=\eps\xi+L(t)V$, $U=QV$, $QQ^T=2I$. Then
\begin{equation}\label{eq:regression}
 \Cov(D_\eps\mid U)=\sigma^2H(t)H(t)^T,\qquad
 H(t)=L(t)\left(I-\tfrac12Q^TQ\right).
\end{equation}
Velocity maps depend on the prescribed order and normals, not on the
contact times inside their ordered domains.

\begin{lemma}[A nonzero contact-time slope]\label{lem:slope}
Let $e_*$ be the earliest edge on the A--B spatial path.
The derivative of the residual coefficient $H$ with respect to its
contact time has operator norm at least one.
\end{lemma}
\begin{proof}
Before $e_*$ the sum of the already-contacted path incidences is zero;
just afterward it is $e_p-e_q$. Differentiating \eqref{eq:position} gives
the negative relative velocity immediately after $e_*$.
Remove this edge from the final tree. The two sides have independent
current Gaussian coordinates, and no later contact crosses the cut.
If $K_A,K_B$ transfer the current $p,q$ velocities to the final A,B
velocities, conditioning on both final velocities gives covariance
\[
 \sigma^2(2I-K_A^TK_A-K_B^TK_B).
\]
At least one transfer is off-diagonal, because the path has at least two
edges. It has rank at most one by \Cref{lem:rankone} and norm at most one
by orthogonality. On its kernel the displayed matrix has a direction of
variance at least $\sigma^2$. Conditioning only on the final difference
increases covariance. Equation \eqref{eq:regression} proves the claim.
\end{proof}

\begin{proposition}[Time-averaged nonneighbor event]\label{prop:nonneighbor}
Let $a_i$ count all free merger times in the full schedule, including
post-return mergers, and let $i_*$ contain $e_*$. The Gaussian prefix
return event $G$ satisfies
\[
 \int_{\prod_i\mathcal S_{a_i}(\tau)}\Pp_\gamma(G)\dd t
 \leq C\prod_i\frac{\tau^{a_i}}{a_i!}
 \min\left(1,\frac{a_{i_*}\eps}{\sigma\tau}
                  [1+\logp(\sigma\tau/\eps)]\right).
\]
This is uniform in all free normals and prescribed forest shapes.
\end{proposition}
\begin{proof}
Vary $t_{e_*}$ between its nearest full-schedule neighbors in its slab,
using slab endpoints if needed. Hold all other parameters fixed.
The slope lemma supplies unit vectors $w,z$ with
$w^TH(t)z=\alpha t_{e_*}+\beta$, $|\alpha|\geq1$.
The largest conditional variance is at least
$\sigma^2|\alpha t_{e_*}+\beta|^2$. For an interval $J$ of length $s$,
\[
 \int_J\min(1,\rho/|\alpha t+\beta|)\dd t
       \leq2\rho[1+\logp(s/(2\rho))],\qquad \rho=\eps/\sigma.
\]
The even decreasing integrand has maximal interval integral when centered
at its singularity. Apply \Cref{lem:smallball}, use $s\leq\tau$, and
remove the selected time from its ordered simplex. The remaining volume
is $\tau^{a_{i_*}-1}/(a_{i_*}-1)!$; every other slab factorial is unchanged.
Also use the trivial probability bound.

For fixed birth and ending interfaces, allocate all birth Gaussian
coordinates in advance, leave unborn coordinates unused, and freely
extend terminated lines only for \eqref{eq:position}. The auxiliary
extensions introduce no contact or random coordinate and cancel in the
spatial telescoping sum. Varying the selected contact within its own slab
retains the original interface assignments.
\end{proof}

An untruncated inverse covariance moment at exponent one is not claimed:
two O contacts on the path $1$--$2$--$3$ give covariance proportional to
the square of their time gap, whose inverse square root is not integrable.
The truncated probability estimate legitimately retains one logarithm.
